% TOP_PROBLEM: {"id": 3344, "title": "Lonely runner conjecture", "category_id": 1, "category": {"id": 1, "name": "number_theory", "display_name": "Number Theory", "description": "Properties of integers, prime numbers, Diophantine equations.", "slug": "number-theory", "order_index": 1, "created_at": "2026-07-31T15:26:25.670Z"}, "status": "open", "problem_number": "OPG-416", "set_id": 12}
% FIRST_POSTED: 2026-10-01
% ENTRY_KIND: statement_only
% UnsolvedMath ID 3344; source catalogue code OPG-416
% !TeX program = lualatex
% Definition and sources only; this file is not an LLM solution attempt.
\documentclass[11pt,a4paper]{article}
\usepackage[margin=25mm]{geometry}
\usepackage{fontspec}
\setmainfont{lmroman10-regular.otf}[BoldFont=lmroman10-bold.otf,
 ItalicFont=lmroman10-italic.otf,BoldItalicFont=lmroman10-bolditalic.otf]
\usepackage{amsmath,amssymb,mathtools}
\usepackage{unicode-math}
\setmathfont{latinmodern-math.otf}
\AtBeginDocument{\let\setminus\smallsetminus}
\usepackage{xurl,hyperref}
\hypersetup{colorlinks=true,urlcolor=blue}
\setcounter{secnumdepth}{0}
\setlength{\parindent}{0pt}
\setlength{\parskip}{5pt}
\setlength{\emergencystretch}{3em}
\begin{document}
\section*{Lonely runner conjecture}\label{3344}
{\small UnsolvedMath ID 3344 · OPG-416 · Number Theory · OpenGarden}
\subsection{Definitions and mathematical statement}
Identify the circular track of circumference one with $\mathbb R/\mathbb Z$.
For $x\in\mathbb R$, let $\|x\|_{\mathbb R/\mathbb Z}
=\min_{m\in\mathbb Z}|x-m|$ be its distance to the nearest integer.
This takes values between zero and $1/2$ and gives the shorter-arc
distance between two points on the track.

Let $k\ge2$ and let $v_1,\ldots,v_k$ be pairwise distinct real speeds.
At time $t\ge0$, runner $i$ occupies $v_it+\mathbb Z$, so all runners
start at zero. The conjecture asserts
\[
 \forall i\in\{1,\ldots,k\}\ \exists t\ge0\quad
 \min_{j\ne i}\|(v_i-v_j)t\|_{\mathbb R/\mathbb Z}\ge\frac1k.
\]
The time can depend on the designated runner and on all the speeds.
The claim does not require all runners to be lonely simultaneously.
Subtracting $v_i$ from every speed makes the designated runner stationary
without changing the displayed distances; negative relative speeds are
therefore allowed. The denominator uses the total number $k$ of runners,
rather than the number $k-1$ of nonzero relative speeds. Equality at
distance $1/k$ is permitted. Establish the assertion for all $k$ and
all such speed vectors, or give a counterexample.
\subsection{Short English statement}
Does each runner eventually stay at least one kth of the track away from every other runner, when there are k distinct speeds?
\subsection{Sources}
\begin{itemize}
\item[S1] ULAM.AI and the UnsolvedMath Contributors, supplied problems.json, record 3344 (OPG-416), OpenGarden collection. Catalogue identity and source transcription; CC BY 4.0.
\url{https://huggingface.co/datasets/ulamai/UnsolvedMath}.
\item[S2] Open Problem Garden, Lonely runner conjecture. Original problem and discussion; the mathematical formulation below is expanded from the supplied source record.
\url{http://www.openproblemgarden.org/op/lonely_runner_conjecture}.
\end{itemize}
\end{document}
